\chapter{Solid State Physics: Final Exam, June 25, 2015}
\label{q:physics-2015-06-final}
\textbf{Printed total: 100\%.}
\begin{sourceNote}
Reconstructed from one complete exam page. Two identical source photographs are represented by this single exam. Handwritten answers are excluded.
\end{sourceNote}

\section*{1. Bloch sums (15\%)}
Express a Bloch sum of wave vector $\vct k$. Show that the Bloch sum has the property of the Bloch function. What is the role the Bloch sums play in the tight-binding method?

\section*{2. Tight-binding model of silicon (15\%)}
The electronic configuration of the silicon atom is
$1s^2 2s^2 2p^6 3s^2 3p^2$.
If five orbitals for each atom are chosen to construct the Bloch sums in the tight-binding method of band structure calculation for electrical and chemical properties of crystalline silicon, what is your suggestion on the five orbitals? How many bands will you have? How many bands will be fully filled by valence electrons for the silicon crystal in the ground state?

\section*{3. Plane-wave expansion (10\%)}
Describe the difficulty in the band structure calculation when plane waves are used as a basis to expand the crystal wavefunction. How to overcome the difficulty?

\section*{4. Pseudopotential (10\%)}
Explain the term ``pseudopotential''. Why can the pseudo wavefunctions for a crystal be expanded with a reasonable number of plane waves?

\section*{5. Adiabatic energy surface (5\%)}
Explain the term ``adiabatic energy surface''.

\section*{6. Normal modes from an energy surface (15\%)}
Consider an electronic-nuclear crystal system on a nondegenerate ground adiabatic surface $E_0(\vct R)$ which well separates in energy from all the other adiabatic surfaces. Describe how to obtain the normal modes of the lattice dynamics with the information of the adiabatic surface $E_0(\vct R)$ and the atomic masses.

\section*{7. Dynamical matrix (5\%)}
Express the dynamical matrix elements, given the atomic masses $M_\nu$ and the force constants
$D_{\nu\alpha,\nu'\alpha'}(\vct t_n)$,
where $\nu,\nu'$ are indices for basis atoms, $\alpha,\alpha'$ are indices for Cartesian coordinate components $x,y,z$, and $\vct t_n$ stands for translation vectors.

\section*{8. Lattice dynamics of BaTiO$_3$ (20\%)}
Consider the BaTiO$_3$ crystal with $N$ unit cells, which has a Ba atom, a Ti atom, and three O atoms in each unit cell.
\begin{enumerate}[label=(\alph*)]
\item What is the dimension of the dynamical matrix?
\item How many normal modes are there of the lattice dynamics?
\item How many acoustic branches are there?
\item How many optical branches are there?
\end{enumerate}

\section*{9. Optical-mode polarization vectors (5\%)}
Consider a crystal with two atoms in each unit cell. The two atoms have masses $M_1$ and $M_2$, respectively. Find the value of
\[
\frac{\abs{\vct A_1}^2}{\vct A_1^*\cdot\vct A_2},
\]
where $\vct A_1$ and $\vct A_2$ are the polarization vectors of the optical phonon modes in the long-wavelength limit $(q\to0)$.

Hint: The atomic displacements of mode $(\vct q,p)$ may be expressed by
\[
\vct u_{n\nu}^{(\vct q,p)}(t)
=\frac{\vct A_\nu(\vct q,p)}{\sqrt{M_\nu}}
w_{\vct q n}\ee^{-\ii\omega(\vct q,p)t}.
\]
